\section{Итог: одна физика, разные модели}
\label{sec:summary}

Хроматичность кольца с секторными диполями и перпендикулярными торцами складывается из пяти членов: A (кинематический, $-\oint\gamma/4\pi$), C (кривизна тела, $+\int D_1\gamma/4\pi\rho$), E (перекрёстный, $-\int D_1'\alpha_x/2\pi\rho$, только $x$), S (секступоли, $\pm\oint K_2D_1\beta/4\pi$) и F (жёсткий край, $\sum[\beta_yD_1']^{\text{вход}}_{\text{выход}}/4\pi\rho$, только $y$). Члены A, C, E, S следуют из гамильтониана тела; F --- из продольного поля $B_s$, которого требует $\operatorname{rot}\mathbf{B}=0$ на границе поля. Разные программы дают разные $\xi_y$ не из-за метода вычисления (все методы главы~\ref{sec:codes} эквивалентны определению главы~\ref{sec:def}), а из-за того, есть ли член F в их модели диполя. Горизонтальная хроматичность от модели края в первом порядке не зависит.

\begin{center}
\small
\renewcommand{\arraystretch}{1.2}
\begin{tabular}{@{}lcccccc@{}}
\hline
Модель & A & C & E & S & F & конечн. край\\
\hline
BELL (Тейлор / Ли) & + & + & + & + & --- & ---\\
COSY \texttt{FR 0} & + & + & + & + & --- & ---\\
COSY \texttt{FR 1}--\texttt{FR 3} & + & + & + & + & + & Энге\\
MAD-X \texttt{TWISS}, \texttt{SBEND} & + & + & + & + & + & \texttt{FINT}, \texttt{HGAP}\\
PTC, \texttt{exact} & + & + & + & + & + & \texttt{fringe=true}\\
OptiM & + & + & + & + & + & ---\\
Python, тело & + & + & + & + & --- & ---\\
Python, тело + \eqref{eq:kick} & + & + & + & + & + & ---\\
\hline
\end{tabular}
\end{center}

Замечания к таблице.
\begin{itemize}
\item Строки <<Python>> --- независимая проверка интегрированием точного гамильтониана \eqref{eq:H} (глава~\ref{sec:codes}).
\item <<Конечный край>> означает поправки $O(g)$ раздела~\ref{sec:edge}; при $g\to0$ строки COSY \texttt{FR 1}--\texttt{FR 3} переходят в строку жёсткого края (MAD-X, PTC, OptiM), а не в строку \texttt{FR 0}.
\item Физически правильная модель диполя с перпендикулярными торцами и коротким спадом поля --- с членом F. Модель <<только тело>> соответствует полю, не удовлетворяющему уравнениям Максвелла на границе.
\item Для COSY с энергией-параметром (\texttt{OV 3 2 1}) и с энергией-координатой (\texttt{OV 3 3 0}, метод неподвижной точки) результат один и тот же; переменная COSY --- $\delta_K$, перевод \eqref{eq:xi-conv}. Для MAD-X \texttt{DQ} относится к $p_t$.
\item С ВЧ линейная $\xi$ получается только сдвигом неподвижной точки \eqref{eq:xi-fp}; \texttt{TS} даёт усреднённую вторую хроматичность \eqref{eq:dnu-dJs}.
\item В $\Lambda$ и равновесный сдвиг \eqref{eq:Lambda} входит полная хроматичность, с членом F.
\item Приближённые модели --- PTC без \texttt{exact}, тонкие диполи после \texttt{MAKETHIN} --- теряют члены с $1/\rho$ (C, E, F) и в таблицу не включены.
\end{itemize}
